\documentclass[11pt,a4paper]{article}
\usepackage[margin=1in]{geometry}
\usepackage{amsmath,amssymb,amsthm}
\usepackage{algorithm}
\usepackage{algpseudocode}
\usepackage{booktabs}
\usepackage{hyperref}

\theoremstyle{definition}
\newtheorem{definition}{Definition}[section]
\newtheorem{theorem}{Theorem}[section]
\newtheorem{lemma}[theorem]{Lemma}
\newtheorem{remark}{Remark}[section]

\title{\textbf{Mathematical Specification, Min-Max Robust Optimization, and First-Order Adversarial Perturbations (FGSM and PGD)}}
\author{Team Scaibu \\ \small Email: \texttt{jaydeep.9664920749@gmail.com} \\ \small Architecture and Core Engine Team}
\date{October 2026}

\begin{document}
\maketitle

\begin{abstract}
This document provides the formal mathematical specification, saddle-point formulation, and projection mechanics for adversarial perturbation algorithms: Fast Gradient Sign Method (FGSM) and Projected Gradient Descent (PGD). Adversarial training reframes standard empirical risk minimization into a robust minimax game $\min_\theta \mathbb{E} [\max_{\|\delta\| \le \epsilon} \mathcal{L}(f_\theta(x + \delta), y)]$. We derive the exact closed-form solution for single-step first-order maximization, formalize multi-step $L_\infty$ ball projection, analyze computational bounds, trace a worked numerical example, and document gradient masking safeguards.
\end{abstract}

\section{Notation and Mathematical Preliminaries}

Let $x \in \mathbb{R}^{B \times D}$ denote a batch of clean input vectors, $y \in \{1, \dots, K\}^B$ the true ground-truth targets, and $\mathcal{L}(\theta; x, y)$ the scalar classification loss.

\begin{itemize}
    \item $B \in \mathbb{N}_{\ge 1}$: Batch size.
    \item $D \in \mathbb{N}_{\ge 1}$: Feature dimensionality.
    \item $\epsilon > 0$: Adversarial perturbation budget bounding the $L_\infty$ ball $\mathcal{S}_\epsilon = \{\delta \in \mathbb{R}^D \mid \|\delta\|_\infty \le \epsilon\}$.
    \item $\alpha > 0$: Step size for iterative gradient ascent in PGD.
    \item $K \in \mathbb{N}_{\ge 1}$: Number of projected gradient ascent iterations ($K = 1$ for FGSM).
    \item $\nabla_x \mathcal{L} \in \mathbb{R}^{B \times D}$: Gradient of the loss with respect to input activations.
    \item $\Pi_{[a, b]}(z) = \max(a, \min(b, z))$: Coordinate-wise projection operator.
    \item $x_{\text{adv}} \in \mathbb{R}^{B \times D}$: Adversarially perturbed input tensor.
\end{itemize}

\begin{table}[h]
\centering
\caption{Summary of Mathematical Symbols for Adversarial Attacks}
\begin{tabular}{lll}
\toprule
\textbf{Symbol} & \textbf{Type / Domain} & \textbf{Description} \\
\midrule
$x$ & $\mathbb{R}^{B \times D}$ & Clean input activation matrix \\
$\nabla_x \mathcal{L}$ & $\mathbb{R}^{B \times D}$ & Loss gradient tensor \\
$\epsilon$ & $\mathbb{R}_{> 0}$ & $L_\infty$ perturbation radius \\
$\alpha$ & $\mathbb{R}_{> 0}$ & Ascent step size \\
$K$ & $\mathbb{N}_{\ge 1}$ & Iterative projection steps \\
$\delta$ & $[-\epsilon, \epsilon]^{B \times D}$ & Adversarial displacement tensor \\
$x_{\text{adv}}$ & $\mathbb{R}^{B \times D}$ & Final perturbed input tensor \\
\bottomrule
\end{tabular}
\end{table}

\section{Problem Definition and Core Concepts}

\subsection{Intuitive Meaning}
Deep neural networks are susceptible to adversarial examples: minute, imperceptible input perturbations designed in the direction of steepest loss ascent that flip predictions with high confidence. Adversarial training formulates optimization as a game against an adversary: finding the worst-case perturbation within an $\epsilon$-ball around each sample, and training parameters $\theta$ to minimize loss on these worst-case inputs.

\subsection{Formal Minimax Optimization Problem}
The robust optimization objective is:
\begin{equation}
\min_\theta \mathbb{E}_{(x, y) \sim \mathcal{D}} \left[ \max_{\|\delta\|_\infty \le \epsilon} \mathcal{L}(f_\theta(x + \delta), y) \right]
\end{equation}

\subsubsection{Fast Gradient Sign Method (FGSM)}
Linearizing $\mathcal{L}(x + \delta)$ around $x$ yields the closed-form single-step maximizer:
\begin{equation}
x_{\text{adv}} = \Pi_{[0, 1]} \left( x + \epsilon \cdot \text{sign}(\nabla_x \mathcal{L}(x, y)) \right)
\end{equation}

\subsubsection{Projected Gradient Descent (PGD)}
Iteratively maximizes loss by stepping along the sign of the local gradient and projecting onto the intersection of the $\epsilon$-ball and the valid pixel domain $[0, 1]$:
\begin{equation}
x^{(t+1)} = \Pi_{[0, 1]} \left( \Pi_{x + [-\epsilon, \epsilon]} \left( x^{(t)} + \alpha \cdot \text{sign}(\nabla_{x^{(t)}} \mathcal{L}(x^{(t)}, y)) \right) \right)
\end{equation}

\section{Algorithm Specification}

\begin{algorithm}[H]
\caption{Iterative Projected Gradient Descent (PGD / FGSM)}
\label{alg:pgd_fgsm}
\begin{algorithmic}[1]
\Require Clean input $X \in \mathbb{R}^{B \times D}$, gradient $G \in \mathbb{R}^{B \times D}$, budget $\epsilon > 0$, step size $\alpha > 0$, iterations $K \in \mathbb{N}_{\ge 1}$, bounds $[c_{\text{min}}, c_{\text{max}}]$.
\Ensure Adversarial tensor $X_{\text{adv}}$, perturbation $\Delta$, max norm $\|\Delta\|_\infty$.
\State $X_{\text{adv}} \gets X$
\State $\alpha_{\text{eff}} \gets \mathbf{if } K = 1 \mathbf{ then } \epsilon \mathbf{ else } \alpha$
\For{$k \gets 1 \textbf{ to } K$}
    \For{$i \gets 1 \textbf{ to } B$}
        \For{$j \gets 1 \textbf{ to } D$}
            \State $s \gets \text{sign}(G_{ij})$
            \State $v \gets X_{\text{adv}}[i][j] + \alpha_{\text{eff}} \cdot s$
            \State $v \gets \max(X[i][j] - \epsilon, \min(X[i][j] + \epsilon, v))$ \Comment{$L_\infty$ Projection}
            \State $X_{\text{adv}}[i][j] \gets \max(c_{\text{min}}, \min(c_{\text{max}}, v))$ \Comment{Domain Clamping}
        \EndFor
    \EndFor
\EndFor
\State $\Delta \gets X_{\text{adv}} - X$
\State $\|\Delta\|_\infty \gets \max_{i, j} |\Delta_{ij}|$
\State \Return $\{X_{\text{adv}}, \Delta, \|\Delta\|_\infty\}$
\end{algorithmic}
\end{algorithm}

\section{Theoretical Guarantees and Optimality}

\begin{theorem}[Dual Norm Optimality of Sign Perturbation]
Let $g = \nabla_x \mathcal{L}(x) \in \mathbb{R}^D$. Under a first-order Taylor approximation $\mathcal{L}(x + \delta) \approx \mathcal{L}(x) + g^\top \delta$, the optimal perturbation $\delta^*$ maximizing loss within the $L_\infty$ constraint $\|\delta\|_\infty \le \epsilon$ is uniquely given by:
\begin{equation}
\delta^* = \epsilon \cdot \text{sign}(g)
\end{equation}
and the inner product achieves the dual $L_1$ norm bound $g^\top \delta^* = \epsilon \|g\|_1$.
\end{theorem}

\begin{proof}
By Hölder's inequality, $g^\top \delta \le \|g\|_1 \|\delta\|_\infty \le \epsilon \|g\|_1$.
Evaluating $\delta^* = \epsilon \cdot \text{sign}(g)$:
\begin{equation}
g^\top \delta^* = \sum_{j=1}^D g_j (\epsilon \cdot \text{sign}(g_j)) = \epsilon \sum_{j=1}^D g_j \text{sign}(g_j) = \epsilon \sum_{j=1}^D |g_j| = \epsilon \|g\|_1
\end{equation}
Because the upper bound is achieved, $\delta^* = \epsilon \text{sign}(g)$ is the exact global maximizer.
\end{proof}

\subsection{Limitations}
\begin{itemize}
    \item \textbf{Obfuscated Gradients}: Defenses relying on non-differentiable operations (shattered gradients) give a false sense of security without conferring true robust accuracy against adaptive attacks.
\end{itemize}

\section{Complexity Analysis}

\subsection{Time Complexity}
\begin{itemize}
    \item \textbf{Sign Computation and Projection}: Updating each element takes $\mathcal{O}(1)$ operations per step.
    \item \textbf{Total Time Complexity}: $\Theta(K \cdot B \cdot D)$ FLOPs.
\end{itemize}

\subsection{Space Complexity}
\begin{itemize}
    \item \textbf{Storage}: $\mathcal{O}(B \cdot D)$ auxiliary space for $X_{\text{adv}}$ and perturbation $\Delta$.
\end{itemize}

\section{Exactness and Error Analysis}

The projection and sign calculations are exact. Clamping prevents pixel values from exceeding legal display limits $[0, 1]$.

\section{Worked Numerical Example}

Let $B = 1, D = 3, \epsilon = 0.1, \alpha = 0.05, K = 2$, bounds $[0.0, 1.0]$.
\begin{equation}
X = [[0.50, 0.95, 0.02]], \quad G = [[1.2, -0.8, 0.5]]
\end{equation}

\begin{enumerate}
    \item $\text{sign}(G) = [1.0, -1.0, 1.0]$.
    \item \textbf{Step 1} ($\alpha = 0.05$):
    \begin{align}
    x^{(1)} &= [0.50 + 0.05, 0.95 - 0.05, 0.02 + 0.05] = [0.55, 0.90, 0.07]
    \end{align}
    All within $[x - 0.1, x + 0.1]$ and $[0, 1]$.
    \item \textbf{Step 2} ($\alpha = 0.05$):
    \begin{align}
    x^{(2)} &= [0.55 + 0.05, 0.90 - 0.05, 0.07 + 0.05] = [0.60, 0.85, 0.12]
    \end{align}
    Bounds check:
    \begin{align}
    x_1^{(2)} &= 0.60 \le 0.50 + 0.10 = 0.60 \quad (\text{exact upper bound}) \\
    x_2^{(2)} &= 0.85 \ge 0.95 - 0.10 = 0.85 \quad (\text{exact lower bound}) \\
    x_3^{(2)} &= 0.12 \le 0.02 + 0.10 = 0.12 \quad (\text{exact upper bound})
    \end{align}
    Perturbation $\Delta = [0.10, -0.10, 0.10]$, $\|\Delta\|_\infty = 0.10$.
\end{enumerate}

\section{Numerical Considerations and Edge Cases}

\begin{itemize}
    \item \textbf{Zero Gradient Elements}: If $G_{ij} = 0$, $\text{sign}(0) = 0$, producing no displacement along flat directions.
\end{itemize}

\begin{thebibliography}{9}
\bibitem{goodfellow2014explaining} I.~J.~Goodfellow, J.~Shlens, and C.~Szegedy, ``Explaining and Harnessing Adversarial Examples,'' in \emph{International Conference on Learning Representations (ICLR)}, 2015.
\bibitem{madry2018towards} A.~Madry, A.~Makelov, L.~Schmidt, D.~Tsipras, and A.~Vladu, ``Towards Deep Learning Models Resistant to Adversarial Attacks,'' in \emph{International Conference on Learning Representations (ICLR)}, 2018.
\end{thebibliography}

\end{document}
